\documentclass[11pt]{article}
\usepackage{amsmath,amssymb,amsthm}
\newtheorem{theorem}{Theorem}
\newtheorem{proposition}[theorem]{Proposition}
\title{Denominator Orders Do Not Determine Full-Cube Sign Classes}
\author{}
\date{November 23, 2026}
\begin{document}
\maketitle

\paragraph{Status and scope.}
The results below are complete written proofs with exact rational
certificates, not Lean formalizations. They audit a proposed shortcut
for the global lower bound. They do not prove either requested global
asymptotic alternative, and they do not alter atom admissibility.

For fixed positive affine denominators $B,C$ on the Boolean cube,
let $\mathcal S(B,C)$ denote the set of strict sign patterns of
$c+A/B+D/C$, with arbitrary real affine numerators and real $c$.

\begin{theorem}
There are two pairs of strictly positive-slope affine denominators on
the four-dimensional Boolean cube such that:
\begin{enumerate}
\item their induced strict orders on all sixteen vertices coincide;
\item even the combined order of all thirty-two denominator values
      coincides between the two pairs;
\item their strict sign classes $\mathcal S(B,C)$ differ.
\end{enumerate}
The distinction persists in open neighborhoods of the two pairs within
the same strict order region. In particular it is not caused merely by
a rank-deficient exceptional choice of denominators.
\end{theorem}
\begin{proof}
Use lexicographic vertex order $0000,0001,\ldots,1111$, indexed from
zero. Define
\[
 B(x)=1+x_1+3x_2+9x_3+27x_4,
 \qquad C_v(x)=100B(x)+v\cdot x,
\]
and take
\[
 v=(8,-9,1,4),\qquad v'=(2,-8,-6,-1).
\]
All slopes and constants of these denominators are strictly positive.
Distinct vertices have distinct integer $B$ values. Since both
$\sum_i|v_i|$ and $\sum_i|v'_i|$ are less than $100$, perturbation
by $v\cdot x$ or $v'\cdot x$ cannot reverse any vertex order of $100B$.
Moreover $B(x)\leq41<100\leq C_v(y),C_{v'}(y)$ for all vertices.
This proves both order assertions.

For either $v$, the sign class is precisely the strict sign class of
the nine features
\[
 1,\quad x_i/B(x),\quad x_i/C_v(x)\quad(1\leq i\leq4).
 \tag{1}
\]
Indeed a constant numerator divided by an affine denominator is in
the span of its four displayed coordinate features and $1$, using
the positive nonzero constant coefficient of that denominator.
Multiplication of each row by $B(x)C_v(x)>0$ gives the integer
feature row
\[
 M_v(x)=\bigl(B(x)C_v(x),\ (x_iC_v(x))_{i=1}^4,
                                  \ (x_iB(x))_{i=1}^4\bigr).
 \tag{2}
\]

Restrict to the ten vertex indices
\[
                 J=(0,1,3,6,8,9,10,12,13,14).
\]
The two $10\times9$ matrices in (2) have rank nine and primitive
integer left-kernel vectors respectively
\[
 \begin{split}
 \lambda={}&(5759190,-584406,319176,-92718,334041,33920,\\
             &\hspace{28mm}-259736,1076448,111568,-553363),\\
 \lambda'={}&(143600853,-7864640,2272334,-574407,-19849030,\\
             &\hspace{28mm}5150082,-10647419,-3965287,-413088,2893753).
 \end{split}
 \tag{3}
\]
These claims have exact certificates: the signed $9\times9$ cofactors
of $M_v|_J$ are $5600\lambda$, and those of $M_{v'}|_J$ are
$400\lambda'$. In particular every cofactor is nonzero. The companion
checker verifies the cofactors and the kernel identities over the
rationals.

Set $\tau_j=\operatorname{sgn}\lambda'_j$. No member of
$\mathcal S(B,C_{v'})$ can restrict to $\tau$ on $J$.
If it did, its cleared numerator vector $z$ would satisfy both
$\lambda'\cdot z=0$ and $\lambda'\cdot z>0$.

Here is an explicit member of $\mathcal S(B,C_v)$ with that
restriction. Take the score
\[
 \begin{split}
 F(x)={}&726397+
 \frac{413029720x_1+1022839140x_2+2018070760x_3+6209249130x_4}{B(x)}\\
 &-\frac{43229541511x_1+100481986842x_2+
              202681766582x_3+623851360053x_4}{C_v(x)}.
 \end{split}
 \tag{4}
\]
Its sixteen strict signs in the stipulated lexicographic order are
\[
                  (+,-,+,+,-,-,-,-,-,+,-,-,-,-,+,-).
 \tag{5}
\]
Exact substitution checks (5), and its restriction to $J$ equals
$\tau$. Thus the two sign classes differ.

Finally all the margins in (4) are strict, all denominator order
inequalities are strict, and all cofactors in (3) are nonzero.
Continuity preserves these facts in sufficiently small neighborhoods
of the respective denominator pairs. On the second neighborhood the
cofactor signs remain $\tau$, so the same strict sign obstruction
persists. On the first neighborhood the fixed score coefficients in
(4) retain the signs in (5). This proves the robustness assertion.
\end{proof}

\begin{proposition}[Quadratic products inside one order region]
Fix any strictly positive-slope affine $B$ with distinct Boolean
vertex values. For any real affine $L,V$ and sufficiently small
$\epsilon>0$, the denominator $C=B+\epsilon V$ has strictly
positive constant and slopes and the same vertex order as $B$.
The exact two-head identity
\[
 \frac{L}{B}-\frac{L}{C}
                  =\frac{\epsilon LV}{BC}
 \tag{6}
\]
therefore realizes the strict sign of any nowhere-zero product $LV$
without leaving that fixed denominator order region.
\end{proposition}
\begin{proof}
There are finitely many strict coefficient and order inequalities.
They all hold at $\epsilon=0$, so they persist for small positive
$\epsilon$. Identity (6) follows by cross multiplication, and $BC$
is positive on the cube.
\end{proof}

The proposition illustrates why arbitrarily small denominator
perturbations cannot be discarded in a sign-count argument: numerator
cancellation can expose their exact directional information. The
theorem goes further by distinguishing generic sign classes within
one strict order region. Neither statement bounds the number of such
classes asymptotically, and neither disproves the desired global lower
bound.

\end{document}
